\documentclass[a4paper,10pt,fleqn]{article}
\usepackage[margin=1in]{geometry}
\usepackage{fuzz}
\begin{document}
\title{Vox Panel Refresh: A Z Specification}
\maketitle

\section{Scope}

The panel polls durable Vox settings only while it has a confirmed Lux
scene. Reading settings, publishing a snapshot for click routing, and
sending a scene to Lux are distinct operations. A successful push
acknowledges the Hub, not the display: even after a push succeeds, an
event from an older widget may arrive. A failed push does not establish
whether Lux kept the old scene, applied the new one, or lost both.

A snapshot names a settings object, not just its value. Its target is
the provider and optional model; its revision names the target and
\emph{ordered} voice roster. A roster key denotes a whole sequence:
swapping distinct voices changes that key, even when membership stays
the same. Revisions are modeled as exact identities:
digest collisions are outside the model. The revision need not change
on a voice-only or notice-only update. Notices themselves are abstracted
to whether a scene refresh was requested. The saved target abstracts
config fields that the write precondition examines, not the whole file.

\begin{zed}
  [SNAPSHOT, TARGET, ROSTER, REVISION]
\also
  SceneState ::= neverShown | installed | disarmed
\also
  Work ::= idle | pollPushDue | foregroundPushDue | correctionDue
\also
  Outcome ::= accepted | failed
\also
  RenderChange ::= identical | patchable | differentTree
\also
  Plan ::= noPush | patch | fullInstall
\also
  NoticeChange ::= noticeSame | noticeChanged
\also
  RecoveryRead ::= diskSnapshot | keepSnapshot
\end{zed}

\begin{axdef}
  targetOf : SNAPSHOT \fun TARGET \\
  rosterOf : SNAPSHOT \fun ROSTER \\
  revisionOf : SNAPSHOT \fun REVISION
\where
  \forall a, b : SNAPSHOT @ \\
  \quad~ revisionOf(a) = revisionOf(b)
    \iff (targetOf(a), rosterOf(a)) = (targetOf(b), rosterOf(b))
\end{axdef}

\section{State and initialization}

The optional confirmed render is the local patch base, not the
display's observed state. An in-flight menu acknowledgement has its own
captured render; its result is not yet confirmed. Both inactive modes
have no patch base; we
distinguish them by history even though \texttt{LiveScene} represents
both with \texttt{None}. The optional capture records the reference used
by \texttt{replace\_if\_current}. The audit relation is ghost state: it
records the event revision and the held revision at each successful
roster-indexed write. Its diagonal invariant rules out stale writes.

\begin{schema}{PanelRefresh}
  phase : SceneState \\
  shown, capture, menu : \power SNAPSHOT \\
  held : SNAPSHOT \\
  saved : TARGET \\
  work : Work \\
  writes : REVISION \rel REVISION
\where
  \# shown \leq 1 \\
  \# capture \leq 1 \\
  \# menu \leq 1 \\
  phase = installed \iff shown \neq \emptyset \\
  capture \neq \emptyset \implies work = idle \\
  writes \subseteq \id REVISION
\end{schema}

\begin{schema}{Init}
  PanelRefresh'
\where
  phase' = neverShown \\
  shown' = \emptyset \\
  capture' = \emptyset \\
  menu' = \emptyset \\
  work' = idle \\
  writes' = \emptyset
\end{schema}

\section{Reading and reconciling}

An external edit does not synchronously alter the panel. A poll starts
only with an installed scene, and captures the held settings by
identity. To exercise the compare-and-replace seam, the model allows
another publisher to replace the held snapshot during the read without
going through the runner. Routed reads and control commits are
serialized by the runner's lock; the menu acknowledgement deliberately
precedes that lock and can overlap a poll. A poll does not start a read
while a scene push is in flight.

\begin{schema}{ExternalEdit}
  \Delta PanelRefresh \\
  target? : TARGET
\where
  saved' = target? \\
  phase' = phase \land shown' = shown \land held' = held \\
  capture' = capture \land work' = work \land writes' = writes \\
  menu' = menu
\end{schema}

\begin{schema}{PollInactive}
  \Xi PanelRefresh
\where
  (phase \neq installed \lor menu \neq \emptyset) \\
  work = idle \land capture = \emptyset
\end{schema}

\begin{schema}{StartPoll}
  \Delta PanelRefresh
\where
  phase = installed \land work = idle \land capture = \emptyset \\
  menu = \emptyset \\
  capture' = \{held\} \\
  phase' = phase \land shown' = shown \land held' = held \\
  saved' = saved \land work' = work \land writes' = writes \\
  menu' = menu
\end{schema}

\begin{schema}{ConcurrentPublish}
  \Delta PanelRefresh \\
  next? : SNAPSHOT
\where
  capture \neq \emptyset \land next? \neq held \\
  held' = next? \\
  phase' = phase \land shown' = shown \land saved' = saved \\
  capture' = capture \land work' = work \land writes' = writes \\
  menu' = menu
\end{schema}

Here \emph{next} is the snapshot proposed by the read. An equal-value
read proposes the existing identity; unreadable disk or an unrelated
roster refusal retains it as well. The proposal is not forced to match
the latest saved target: the file can change again after the read. A
changed notice alone can schedule a refresh. An obsolete read swaps
neither settings nor notices and schedules no push.

\begin{schema}{AcceptPoll}
  \Delta PanelRefresh \\
  next? : SNAPSHOT \\
  notice? : NoticeChange
\where
  capture = \{held\} \\
  held' = next? \land capture' = \emptyset \\
  work' = \IF next? \neq held \lor notice? = noticeChanged
    \THEN pollPushDue \ELSE idle \\
  phase' = phase \land shown' = shown \land saved' = saved \\
  writes' = writes \land menu' = menu
\end{schema}

\begin{schema}{DiscardPoll}
  \Delta PanelRefresh
\where
  capture \neq \emptyset \land capture \neq \{held\} \\
  capture' = \emptyset \\
  phase' = phase \land shown' = shown \land held' = held \\
  saved' = saved \land work' = work \land writes' = writes \\
  menu' = menu
\end{schema}

\section{Publishing to Lux}

A queued refresh plans no push for equivalent renders, patches a
compatible tree, or installs the whole scene. Without a patch base,
the only plan is a full install; only a foreground event can queue
such a push. Different revisions imply
different widget IDs and cannot be patched in place. A patch refusal
can fall back to a full install; \emph{outcome} below is the eventual
result, after any fallback. Any final refusal or transport loss
disarms polling. The push lock makes scene writes sequential. A poll
that began while installed rechecks the patch base after a concurrent
menu install completes; only foreground work may arm a disarmed
scene. The menu's follow-up read under the routing lock is abstracted
away here; it is foreground work, not a polling tick.

\begin{schema}{PushRefresh}
  \Delta PanelRefresh \\
  change? : RenderChange \\
  outcome? : Outcome \\
  plan! : Plan
\where
  menu = \emptyset \\
  work = foregroundPushDue \\
  \lor (work = pollPushDue \land phase = installed) \\
  phase \neq installed \implies change? = differentTree \\
  (\exists old : shown @ revisionOf(old) \neq revisionOf(held))
    \implies change? = differentTree \\
  plan! = \IF change? = identical \THEN noPush \\
  \quad~ \ELSE (\IF change? = patchable \THEN patch \ELSE fullInstall) \\
  plan! = noPush \implies outcome? = accepted \\
  phase' = \IF outcome? = accepted \THEN installed \ELSE disarmed \\
  shown' = \IF outcome? = accepted \THEN \{held\} \ELSE \emptyset \\
  work' = idle \land capture' = capture \\
  held' = held \land saved' = saved \land writes' = writes \\
  menu' = menu
\end{schema}

\begin{schema}{SkipDisarmedPollPush}
  \Delta PanelRefresh
\where
  work = pollPushDue \land menu = \emptyset \land phase \neq installed \\
  work' = idle \\
  phase' = phase \land shown' = shown \land held' = held \\
  saved' = saved \land capture' = capture \land writes' = writes \\
  menu' = menu
\end{schema}

\begin{schema}{BeginMenuInstall}
  \Delta PanelRefresh
\where
  menu = \emptyset \\
  menu' = \{held\} \\
  phase' = phase \land shown' = shown \land held' = held \\
  saved' = saved \land capture' = capture \land writes' = writes \\
  work' = work
\end{schema}

\begin{schema}{CompleteMenuInstall}
  \Delta PanelRefresh \\
  outcome? : Outcome
\where
  menu \neq \emptyset \\
  phase' = \IF outcome? = accepted \THEN installed \ELSE disarmed \\
  shown' = \IF outcome? = accepted \THEN menu \ELSE \emptyset \\
  held' = held \land saved' = saved \land writes' = writes \\
  capture' = capture \land work' = work \land menu' = \emptyset
\end{schema}

\begin{schema}{CorrectScene}
  \Delta PanelRefresh \\
  outcome? : Outcome
\where
  work = correctionDue \land menu = \emptyset \\
  phase' = \IF outcome? = accepted \THEN installed \ELSE disarmed \\
  shown' = \IF outcome? = accepted \THEN \{held\} \ELSE \emptyset \\
  held' = held \land saved' = saved \land writes' = writes \\
  capture' = capture \land work' = idle \land menu' = menu
\end{schema}

\section{Roster-indexed events}

The runner checks the event ID against the scene projected from
\emph{held} settings, not the last push acknowledged by Lux. An
unrevisioned event fails the same check as a stale one. A matching ID
reaches the commit path, where the provider and model must still
match disk under the write lock. Only a valid choice that succeeds
through those checks enters the commit operation; a failed choice
does not persist a setting. The next snapshot represents the selected
voice or the cascaded model, voice and roster. For a Voice pick its
target stays unchanged. The operation is independent of $phase$:
queued events can arrive even after disarm. The foreground correction
or re-push that follows a routed event is modeled separately.

\begin{schema}{RejectStaleRevision}
  \Delta PanelRefresh \\
  revision? : \power REVISION
\where
  capture = \emptyset \land work = idle \\
  \# revision? \leq 1 \\
  revision? \neq \{revisionOf(held)\} \\
  work' = correctionDue \\
  phase' = phase \land shown' = shown \land held' = held \\
  capture' = capture \land saved' = saved \land writes' = writes \\
  menu' = menu
\end{schema}

\begin{schema}{RejectStaleTarget}
  \Delta PanelRefresh \\
  revision? : REVISION \\
  recovered? : SNAPSHOT \\
  read? : RecoveryRead
\where
  capture = \emptyset \land work = idle \\
  revision? = revisionOf(held) \land targetOf(held) \neq saved \\
  held' = \IF read? = diskSnapshot \THEN recovered? \ELSE held \\
  phase' = phase \land shown' = shown \land saved' = saved \\
  capture' = capture \land work' = correctionDue \land writes' = writes \\
  menu' = menu
\end{schema}

A target mismatch takes the write-failure recovery path, unlike a
revision mismatch. It does not write, but it rereads the actual
config and roster under the runner's routing lock before the
corrective full install. A roster refusal can still yield a disk
snapshot with no roster; an unreadable file or unavailable daemon
retains the held snapshot. The write-failure notice accompanies either
outcome. The file can change again after the recovery read, so the
recovered target is not required to equal the latest saved target.

\begin{schema}{CommitRosterChoice}
  \Delta PanelRefresh \\
  revision? : REVISION \\
  next? : SNAPSHOT
\where
  capture = \emptyset \land work = idle \\
  revision? = revisionOf(held) \land targetOf(held) = saved \\
  held' = next? \land saved' = targetOf(next?) \\
  writes' = writes \cup \{revision? \mapsto revisionOf(held)\} \\
  phase' = phase \land shown' = shown \\
  capture' = capture \land work' = foregroundPushDue \land menu' = menu
\end{schema}

The provider picker does not use a roster revision. A successful
provider switch persists and publishes its new target together, then
queues a push. This can also leave an earlier Voice or Model event
waiting to be delivered. The operation abstracts the validation and
roster fetch that must already have succeeded.

\begin{schema}{CommitProvider}
  \Delta PanelRefresh \\
  next? : SNAPSHOT
\where
  capture = \emptyset \land work = idle \\
  targetOf(next?) \neq targetOf(held) \\
  held' = next? \land saved' = targetOf(next?) \\
  phase' = phase \land shown' = shown \land writes' = writes \\
  capture' = capture \land work' = foregroundPushDue \land menu' = menu
\end{schema}

\section{Consequences and test correspondence}

\begin{description}
\item[DiscardPoll] The settings test rejects an old resync without
replacing the current snapshot or its notice. The guard compares
identities even when revisions are equal.
\item[RejectStaleRevision] The runner seam test delivers delayed Voice
and Model picks after both external and in-panel provider changes.
\item[RejectStaleTarget] The runner seam test switches the disk provider
before polling, then sends Voice and Model events from the old scene.
\item[PollInactive] The service test requires a foreground install
before polling resumes after either a push refusal or unavailability.
\item[SkipDisarmedPollPush] The runner seam test lets a menu install
fail while a poll is reading and verifies that the poll does not
attempt a replacement install.
\end{description}

This policy prioritizes not pushing blindly after losing the patch
base. If Lux still shows the old render, a transient failure can
therefore leave it stale until a foreground action. Recovery without
foreground action would require a way to re-establish the scene's
presence, not just a successful settings read. It is a liveness
limitation, not a violation of the stale-event safety invariant.

\end{document}
